\documentclass{article}
\usepackage{amsmath,hyperref}
\title{When a failed joint attempt consumes the right to fall back}
\date{}
\begin{document}
\maketitle

\paragraph{What changes.}
The companion model conditions pricing on a timely safety certificate. Here a crew must accept or reject a tentative joint job at tick 4; neither party has yet received the required clearance. The actual shared launch at tick 5 is gated on recognized authority and timely truthful evidence that an old physical use has ended. A failed read or an occupied passage cannot be waived with a price. A transfer is payable only when the shared service is physically performed. A failed attempt may nonetheless consume a staging slot, leaving no time to carry out the stand-alone fallback. The stakeholder law preserves the old accepted work and the new clearance requirement in every history; participation and the probability of optional useful service are separate.

The fixture has a one-tick staging action at 4 for the fallback and a one-tick completion after a failed shared attempt at 5, due by 6. A known valid permission and actual clearance on the fallback path must permit that staging. Rejecting the offer leaves the staging lane free. When the tentative attempt and fallback must use the same lane at 4, accepting prevents the fallback from ever becoming ready, even if the shared launch never starts. Alternatively an already available independent staging lane can be reserved at a prior resource cost of one, which blocks other optional use of it. Now the crew can stage the fallback when its own permission and clearance hold, finish it if shared clearance fails, or retire the staged kit on that independent lane while the shared job proceeds. Teardown uses the reserved lane on $[5,6)$ and is not a second delivered service. The rule observes only its own fallback evidence and later lawfully received clearance. The comparison does not silently reclaim the lane during the unsuccessful branch.

\paragraph{Exact willingness to pay.}
Let $s_i>0$ be the conditional probability of actual lawful performance if type $i$ accepts, $a_i$ the net value of immediate rejection, and $\ell_i$ the expected fallback after failure, conditional on failure and that same type and attempt. Let $b_i\ge0$ be the seller's expected nonrecoverable attempt cost, and $k$ its actual performance cost. If $s_i=1$ the product $(1-s_i)\ell_i$ is zero and does not require defining a conditional distribution on an impossible event. At price $p$, buyer acceptance requires
\[
s_i(v-p)+(1-s_i)\ell_i\ge a_i.
\]
The effective cap is therefore $P_i=v-[a_i-(1-s_i)\ell_i]/s_i$. When $\ell_i=a_i$ we recover the old cap $v-a_i$ even if the service often fails. When failed setup loses the fallback, the cap becomes $v-a_i/s_i$. These are derived reservation values, not observations available to a remote crew before it gets the applicable information. A type with $s_i=0$ is outside the division rule and needs separate allowance for refusing an always-failed attempt.

Pricing on marginal safety and marginal fallback availability is unsound. Two source laws can each have half probability of actual clearance and half probability of a ready fallback after staging. In one law the fallback is ready precisely in the failed shared branch, and in the other it is ready only where the shared attempt succeeds. With $v=18$ and $a_i=6$, the two caps are respectively 12 and 6. Treating readiness and clearance as independent gives cap 9. This comparison concerns what happens after a tentative contract with the same underlying law, not the effect of intentionally violating a safety rule. The invoice is settled from the same history as the actual effect.

\paragraph{Strong matched economic reference.}
At a fee information cell use masses $\pi_i$. For any common price, seller profit is
\[
 \sum_{i:P_i\ge p}\pi_i[s_i(p-k)-b_i].
\]
With acceptance at equality, maximum is among the caps and no attempt; between adjacent caps the selected set is constant and the objective increases. Types receiving equal effective caps have identical preferences over every menu with only attempt probability and fees conditional on actual performance. Differing utility scales $s_i$ do not distinguish their absolute preferred menu choices. We do not assume an independent certification of hidden types, or that arbitrary ties reveal them. At equal caps a mechanism that only has this fee interface cannot demand different allocations backed only by weak truth-telling.

Under that interface restriction, any incentive-compatible direct lottery can collapse equal caps. Order the remaining caps $P_1<\cdots<P_n$. Incentive constraints imply nondecreasing attempt probabilities $x_i\in[0,1]$. Use utilities relative to immediate rejection and normalized by $s_i$. Minimum feasible utilities are $u_1=0$ and $u_{i+1}=u_i+(P_{i+1}-P_i)x_i$; the adjacent lower bounds imply all pairwise constraints by monotonicity. Increasing these utilities above the minimum only reduces seller revenue. After substitution the revenue is linear in the ordered $x_i$:
\[
 \sum_i\pi_i\{s_i[(P_i-k)x_i-u_i]-b_ix_i\}.
\]
Every monotone vector is a convex combination of threshold allocations and no attempt, with weights the adjacent increments and the unused mass. The optimal value is therefore attained by a posted cap or no attempt. If exact acceptance ties cannot be assumed, the corresponding upper limits can be approached with arbitrarily small fee reductions in positive-margin cases. The exact implementation scope does not include menus conditioning fees on independently certified type or failure events; adding those powers changes the optimization.

A numerical weak-truth comparator can overstate this value. Two equally likely types have cap 12. The first has $s=1/2$, $b=6$; the second has $s=1$, $b=0$; $k=4$. An oracle selecting just the second makes four gross. Either type could choose the same zero-rent nominal offer of attempt at 12; no truthful report can be forced from those preferences. The admissible common posted price earns three. The two types may of course be distinguished by another lawful timely test; then three is not a bound on that wider interface.

\paragraph{A paid positive boundary.}
Retain the previous values $v=18,k=4,a=6$ with two equally likely options (none or six), within each type equally likely realized shared clearance and no clearance. A must see enough surplus to cover setup four. Under single-lane staging, accepting loses a valuable fallback if no shared launch is lawful. The seller therefore quotes 18 and only the zero-option type accepts; gross A earns three and declines setup. With a paid independent staging reservation, the effective caps are 18 and 12. Price 12 is optimal. A's gross return is $4.5$, and B's gross return four covers grant overhead two and additional reservation cost one. Both accept the install decision. Conditional on setup, actual shared service frequencies are $1/4$ and $1/2$ in the two configurations respectively. Without setup there is no service. The paid reservation preserves all known old duties and changes useful service without weakening clearance. A factored cached controller with the same reserve and authentic observations makes the same decision; this does not require a new planner.

\paragraph{Comparison and checks.}
This uses the same screening reasoning as \href{https://www.cs.princeton.edu/courses/archive/spring09/cos444/papers/myerson81.pdf}{standard single-object incentive-compatible design} and \href{https://doi.org/10.1016/j.jet.2016.10.008}{bargaining with outside options}. Neither generic equilibrium nor legal permission alone determines which $a_i,\ell_i$ the source supports. The replay checks 2,100 explicit price and same-law type/payoff cases, 120 generic monotone-incentive numerical LP cases and the separate equal-preference oracle control, 128 physical event traces, and the exact factorization counterexample. Lean checks the price cutoff, preference scaling and safety guard implication. The proof does not require belief evolution within the one offer, does not verify network or regulator trust, and is not a general market design result or R5 acceptance.
\end{document}
